\lyseis{Γ1.3}{Κυκλικά διαγράμματα και ιστογράμματα}
\sol{1} α) $\frac{10}{40}=25\%$. β) $\hat{\omega}=\frac{15}{40}\cdot360^\circ=135^\circ$. γ) $\frac{5}{40}\cdot360^\circ=45^\circ$.
\sol{2} $\frac{14}{40}=35\%$ και $\frac{14}{40}\cdot360^\circ=126^\circ$. $\frac{10}{40}=25\%$ και $90^\circ$. $\frac{12}{40}=30\%$ και $108^\circ$. $\frac{4}{40}=10\%$ και $36^\circ$. Άθροισμα γωνιών $126+90+108+36=360^\circ$.
\sol{3} Ο τομέας $90^\circ$ είναι το $\frac{1}{4}$ του κύκλου και έχει $12$ παιδιά, άρα η ομάδα έχει $4\cdot12=48$ παιδιά. Κάθε παιδί αντιστοιχεί σε $7{,}5^\circ$. Μουσική: $120:7{,}5=16$. Θέατρο: $8$. Χορός: $8$. Άλλο: $4$. Έλεγχος: $12+16+8+8+4=48$.
\sol{4} α) $5$ s. β) Η κλάση $40-45$ (συχνότητα $16$). γ) $4+10=14$ παιδιά. δ) $6+4=10$ παιδιά, δηλαδή $\frac{10}{40}=25\%$.
\sol{5} α) Η βάση ισούται με το πλάτος της κλάσης και το ύψος με τη συχνότητα της κλάσης. β) $90^\circ$. γ) $360^\circ$. δ) Ακουμπούν.
\sol{6} α) Κυκλικό διάγραμμα (ποσοστά, λίγες κατηγορίες). β) Ιστόγραμμα (συνεχή, ομαδοποιημένα δεδομένα). γ) Ραβδόγραμμα (πολλές κατηγορίες).
\sol{7} α) Σωστό. β) Λάθος: το ύψος είναι η συχνότητα. γ) Σωστό. δ) Λάθος: ταιριάζει σε λίγες κατηγορίες. ε) Σωστό.
\sol{8} Ελεύθερη κατασκευή. Τα ορθογώνια είναι κολλητά με βάση $2$ cm (το πλάτος της κλάσης) και ύψη $2$, $5$, $8$, $4$, $1$.
\sol{9} $\frac{24}{60}\cdot360^\circ=144^\circ$, $\frac{15}{60}\cdot360^\circ=90^\circ$, $\frac{12}{60}\cdot360^\circ=72^\circ$, $\frac{9}{60}\cdot360^\circ=54^\circ$. Άθροισμα $360^\circ$. Σχετικές συχνότητες $40\%$, $25\%$, $20\%$, $15\%$. Ελεύθερη κατασκευή του κυκλικού διαγράμματος με μοιρογνωμόνιο.
\sol{10} $\frac{6}{40}=15\%$, $\frac{12}{40}=30\%$, $\frac{14}{40}=35\%$, $\frac{6}{40}=15\%$, $\frac{2}{40}=5\%$. Το ιστόγραμμα σχετικών συχνοτήτων έχει κολλητά ορθογώνια με ύψη $15$, $30$, $35$, $15$, $5$ (ελεύθερη κατασκευή).
\sol{11} Κάθε παιδί αντιστοιχεί σε $\frac{360^\circ}{36}=10^\circ$. Ποδόσφαιρο: $150:10=15$ παιδιά. Μπάσκετ: $9$. Κολύμβηση: $7$. Άλλο: $5$. Έλεγχος: $15+9+7+5=36$. Το ποδόσφαιρο προτιμά το $\frac{15}{36}=\frac{5}{12}$ των παιδιών.
\sol{12} α) $6+15=21$ επισκέπτες, δηλαδή $\frac{21}{60}=35\%$. β) $12+6=18$ επισκέπτες, δηλαδή $\frac{18}{60}=30\%$. γ) Στην κλάση $10-20$, που έχει τη μεγαλύτερη συχνότητα ($15$).
\sol{13} Με $9$ χρώματα και μικρούς τομείς είναι δύσκολο να συγκρίνει κανείς τους τομείς και να διακρίνει τους πολύ μικρούς. Καλύτερο είναι ένα ραβδόγραμμα, όπου οι ράβδοι συγκρίνονται εύκολα.
\sol{14} Το $100\%$ αντιστοιχεί σε $360^\circ$, άρα το $1\%$ σε $360^\circ:100=3{,}6^\circ$. Για $5\%$: $5\cdot3{,}6^\circ=18^\circ$. Για $12\%$: $12\cdot3{,}6^\circ=43{,}2^\circ$. Για $37{,}5\%$: $37{,}5\cdot3{,}6^\circ=135^\circ$.
\sol{15} Νέες συχνότητες: $1+6=7$, $6+6=12$, $4+2=6$. Σύνολο $7+12+6=25$. Το ιστόγραμμα έχει τρία πλατύτερα ορθογώνια με ύψη $7$, $12$, $6$. Μένει ίδιο το σύνολο και το γενικό σχήμα (τα περισσότερα παιδιά είναι στη μέση). Χάνονται όμως λεπτομέρειες, π.χ. ότι η πρώτη κλάση είχε μόνο ένα παιδί.
\sol{16} Σωστές γωνίες: $\frac{8}{24}\cdot360^\circ=120^\circ$, $\frac{6}{24}\cdot360^\circ=90^\circ$, $\frac{7}{24}\cdot360^\circ=105^\circ$, $\frac{3}{24}\cdot360^\circ=45^\circ$. Λάθος είναι η γωνία $110^\circ$ και πρέπει να είναι $105^\circ$. Έλεγχος: $120+90+105+45=360$.
\sol{17} Η μέρα έχει $24$ ώρες και ο κύκλος $360^\circ$, άρα η μία ώρα αντιστοιχεί σε $\frac{360^\circ}{24}=15^\circ$. Παράδειγμα: ύπνος $9$ ώρες ($135^\circ$), σχολείο $7$ ώρες ($105^\circ$), φαγητό και μετακινήσεις $3$ ώρες ($45^\circ$), διάβασμα $2$ ώρες ($30^\circ$), ελεύθερος χρόνος $3$ ώρες ($45^\circ$). Άθροισμα ωρών $24$ και γωνιών $360^\circ$.
\sol{18} Ελεύθερη απάντηση. Τα συμπεράσματα πρέπει να στηρίζονται σε αυτό που δείχνει το διάγραμμα (π.χ. ποια κατηγορία ή ποια κλάση έχει τη μεγαλύτερη συχνότητα, ποιο ποσοστό είναι κάτω από κάποιο όριο).
